\documentclass[12pt,a4paper]{article}
\usepackage[utf8]{vietnam}
\usepackage{amsmath,amssymb}
\usepackage{geometry}
\geometry{top=1.5cm,bottom=1.2cm,left=1.4cm,right=1.4cm,headheight=28pt,headsep=12pt,footskip=18pt}
\usepackage{tikz}
\usepackage{xcolor}
\usepackage{enumerate}
\usepackage{graphicx}
\usepackage[version=4]{mhchem}
\usepackage{fancyhdr}
\usepackage{ex_test}

\renewcommand{\baselinestretch}{0.95}
\setlength{\parskip}{0pt}

\definecolor{hfRed}{RGB}{211,47,47}
\definecolor{hfGreen}{RGB}{139,195,144}

\pagestyle{fancy}
\fancyhf{}
\fancyhead[L]{\small\itshape\color{hfRed} Tài liệu lưu hành nội bộ}
\fancyhead[R]{\small\itshape\color{hfRed} Thầy \textbf{TRẦN VĂN THẠNH} - 0777.470.803}
\renewcommand{\headrulewidth}{0.8pt}
\renewcommand{\headrule}{\hbox to\headwidth{\color{hfGreen}\leaders\hrule height \headrulewidth\hfill}}

\fancyfoot[L]{\small\itshape\color{hfRed} CS1: 6/15 Nguyễn Hoàng, P. Kim Long, TP Huế \quad|\quad CS2: 24 Đặng Thái Thân, TP Huế}
\fancyfoot[R]{\small\bfseries\color{hfRed} Trang \thepage}
\renewcommand{\footrulewidth}{0.8pt}
\renewcommand{\footrule}{\hbox to\headwidth{\color{hfGreen}\leaders\hrule height \footrulewidth\hfill}}

\fancypagestyle{firstpage}{
  \fancyhf{}
  \renewcommand{\headrulewidth}{0pt}
  \fancyfoot[L]{\small\itshape\color{hfRed} CS1: 6/15 Nguyễn Hoàng, P. Kim Long, TP Huế \quad|\quad CS2: 24 Đặng Thái Thân, TP Huế}
  \fancyfoot[R]{\small\bfseries\color{hfRed} Trang \thepage}
  \renewcommand{\footrulewidth}{0.8pt}
  \renewcommand{\footrule}{\hbox to\headwidth{\color{hfGreen}\leaders\hrule height \footrulewidth\hfill}}
}

\begin{document}
\thispagestyle{firstpage}

\noindent
\begin{minipage}[t]{0.42\textwidth}
	\centering\fontsize{10pt}{12pt}\selectfont
	\textbf{SỞ GD \& ĐT THÀNH PHỐ HUẾ}\\
	\textbf{TRƯỜNG THPT HOÁ CHÂU}\\
	\textbf{ĐỀ CHÍNH THỨC}\\
	\textit{(Đề gồm có 03 trang)}
\end{minipage}\hfill
\begin{minipage}[t]{0.56\textwidth}
	\centering\small
	\textbf{ĐỀ KIỂM TRA GIỮA HỌC KÌ I – NĂM HỌC 2025-2026}\\
	\textbf{MÔN: HOÁ HỌC – LỚP 12}\\
	\textit{Thời gian làm bài: 45 phút (không kể thời gian phát đề)}\\
	\fbox{\textbf{MÃ ĐỀ: 485}}
\end{minipage}

\smallskip
\noindent\hrulefill
\smallskip

\noindent\textbf{Họ, tên thí sinh:} \dotfill\ \textbf{Lớp:} \dotfill\ \textbf{Số báo danh:} \dotfill

\smallskip
\noindent\textit{(Cho nguyên tử khối của các nguyên tử: H = 1; C = 12; O = 16; N = 14; Cl = 35,5)}
\smallskip

\noindent
\textbf{A. TRẮC NGHIỆM KHÁCH QUAN}

\smallskip
\noindent
\textbf{PHẦN I (3 điểm -- 12 câu): Câu trắc nghiệm nhiều phương án lựa chọn.} \textit{Thí sinh trả lời từ câu 1 đến câu 12. Mỗi câu hỏi thí sinh chỉ chọn một phương án.}

\Opensolutionfile{ans}[ans-hs485]

\begin{ex}
Thuỷ phân $324\text{ gam}$ tinh bột với hiệu suất 75\% khối lượng glucose thu được là
\choice
{$360\text{ gam}$}
{\True $270\text{ gam}$}
{$300\text{ gam}$}
{$250\text{ gam}$}
\loigiai{}
\end{ex}

\begin{ex}
Phát biểu nào sau đây không đúng khi nhận xét về maltose?
\choice
{Là một disaccharide}
{Phân tử có một nhóm -OH hemiacetal}
{Cho được phản ứng thuỷ phân}
{\True Không làm mất màu nước bromine}
\loigiai{}
\end{ex}

\begin{ex}
Khi xà phòng hoá tristearin thu được sản phẩm là
\choice
{$\ce{C15H31COONa}$ và $\ce{C2H5OH}$}
{$\ce{C15H31COOH}$ và $\ce{C3H5(OH)3}$}
{\True $\ce{C17H35COONa}$ và $\ce{C3H5(OH)3}$}
{$\ce{C17H35COOH}$ và $\ce{C3H5(OH)3}$}
\loigiai{}
\end{ex}

\begin{ex}
Công thức cấu tạo dạng mạch vòng của $\alpha$-glucose là
\begin{center}
	\includegraphics[width=14.5cm]{San_Pham/images_crop_485/fig_cau4_glucose_clean.png}
\end{center}
\choice
{\True Cấu trúc A}
{Cấu trúc B}
{Cấu trúc C}
{Cấu trúc D}
\loigiai{}
\end{ex}

\begin{ex}
Đun nóng hỗn hợp các chất trong bình cầu chứa $12\text{ mL}$ acetic acid ($D = 1{,}05\text{ g/mL}$), $11\text{ mL}$ pentyl alcohol ($D = 0{,}81\text{ g/mL}$) và $4\text{ mL}$ dung dịch $\ce{H2SO4}$ đặc cùng một ít đá bọt trong 20 phút. Khối lượng ester thu được sau khi tách khỏi hỗn hợp và làm sạch là $8\text{ gam}$. Hiệu suất phản ứng ester hoá đạt
\choice
{$29{,}30\%$}
{\True $60{,}78\%$}
{$66{,}67\%$}
{$72{,}27\%$}
\loigiai{}
\end{ex}

\begin{ex}
Chất nào sau đây không phải là ester?
\choice
{$\ce{HCOOCH3}$}
{\True $\ce{C2H5COC3H7}$}
{$\ce{CH3COOC2H5}$}
{$\ce{(C15H31COO)3C3H5}$}
\loigiai{}
\end{ex}

\begin{ex}
Cho các phát biểu sau:
\begin{enumerate}[(a)]
	\item Khi làm đậu phụ xảy ra sự đông tụ protein.
	\item Enzyme bị biến tính không thể thực hiện vai trò xúc tác.
	\item Không nên vắt chanh vào sữa khi uống.
	\item Sự thuỷ phân protein xảy ra trong quá trình làm nước mắm hay nấu nước tương.
	\item Mỗi enzyme có một nhiệt độ tối ưu. Tại nhiệt độ tối ưu, enzyme có hoạt tính tối đa làm tốc độ phản ứng xảy ra nhanh nhất.
\end{enumerate}
Số phát biểu đúng là:
\choice
{2}
{3}
{4}
{\True 5}
\loigiai{}
\end{ex}

\begin{ex}
Tên gọi của hợp chất $\ce{C6H5-CH2-CH(NH2)-COOH}$ là
\choice
{Amino phenyl propionic acid}
{phenylalanine}
{2-amino-3-phenylpropionic acid}
{\True 2-amino-3-phenylpropanoic acid}
\loigiai{}
\end{ex}

\begin{ex}
Công thức cấu tạo của glycine là
\choice
{$\ce{C2H5NH2}$}
{\True $\ce{H2NCH2COOH}$}
{$\ce{CH3NH2}$}
{$\ce{H2NCH(CH3)COOH}$}
\loigiai{}
\end{ex}

\begin{ex}
Cho dãy các chất: tinh bột, cellulose, glucose, fructose, saccharose. Số chất trong dãy thuộc loại monosaccharide là
\choice
{1}
{3}
{\True 2}
{4}
\loigiai{}
\end{ex}

\begin{ex}
Chất nào sau đây là thành phần chính của chất giặt rửa tổng hợp?
\begin{center}
	\includegraphics[width=6.5cm]{San_Pham/images_crop_485/fig_cau11_detergent_clean.png}
\end{center}
\choice
{$\ce{(C17H35COO)2Ca}$}
{$\ce{C15H31COONa}$}
{$\ce{C17H35COOK}$}
{\True $\ce{CH3[CH2]11-C6H4-SO3Na}$}
\loigiai{}
\end{ex}

\begin{ex}
Số amine bậc I trong số các chất: $\ce{CH3NH2}$, $\ce{CH3NH3Cl}$, $\ce{(NH2)2CO}$, $\ce{CH3NHCH3}$, $\ce{CH3CH2NH2}$, $\ce{NH2CH2NH2}$, $\ce{(CH3)3N}$, $\ce{C6H5NH2}$ (aniline) là:
\choice
{6}
{7}
{5}
{\True 4}
\loigiai{}
\end{ex}

\Closesolutionfile{ans}

\vspace{0.3cm}
\noindent
\textbf{PHẦN II (2 điểm -- 2 câu): Câu trắc nghiệm đúng sai.} \textit{Thí sinh trả lời từ câu 1 đến câu 2. Trong mỗi ý a), b), c), d) ở mỗi câu, thí sinh chọn đúng hoặc sai.}

\setcounter{ex}{0}

\begin{ex}
Thực hiện phản ứng điều chế isoamyl acetate theo trình tự sau:
\begin{itemize}
	\item \textit{Bước 1:} Cho $2\text{ mL}$ isoamyl alcohol, $2\text{ mL}$ acetic acid nguyên chất và 2 giọt sulfuric acid đặc vào ống nghiệm khô.
	\item \textit{Bước 2:} Lắc đều, đun cách thủy hỗn hợp 5-6 phút trong nồi nước nóng.
	\item \textit{Bước 3:} Để nguội, rồi rót hỗn hợp sản phẩm vào ống nghiệm chứa $2\text{ mL}$ dung dịch $\ce{NaCl}$ bão hòa.
\end{itemize}
\begin{enumerate}[a)]
	\item Tách isoamyl acetate từ hỗn hợp sau bước 3 bằng phương pháp chiết.
	\item Mục đích chính của việc cho dung dịch $\ce{NaCl}$ bão hòa nhằm tránh sự phân hủy sản phẩm.
	\item Ở bước 2 xảy ra phản ứng ester hóa, giải phóng hơi có mùi thơm của chuối chín.
	\item Phản ứng ester hóa giữa isoamyl alcohol với acetic acid là phản ứng thuận nghịch.
\end{enumerate}
\loigiai{}
\end{ex}

\begin{ex}
Em hãy cho biết các phát biểu sau đúng hay sai?
\begin{enumerate}[a)]
	\item Thuỷ phân saccharose và maltose chỉ thu được một monosaccharide duy nhất.
	\item Glucose và Fructose đều có thể mở vòng thành dạng mạch hở.
	\item Glucose và fructose là đồng phân của nhau.
	\item Thuỷ phân 1 tấn tinh bột thu được 10 tấn glucose. Hiệu suất thuỷ phân tinh bột đạt 80\%.
\end{enumerate}
\loigiai{}
\end{ex}

\vspace{0.3cm}
\noindent
\textbf{PHẦN III (2,0 điểm -- 8 câu): Câu trắc nghiệm yêu cầu trả lời ngắn.} \textit{Thí sinh trả lời từ câu 1 đến câu 8.}

\setcounter{ex}{0}

\begin{ex}
Lên men $m\text{ gam}$ glucose với hiệu suất 90\%, lượng khí $\ce{CO2}$ sinh ra hấp thụ hết vào nước vôi trong, thu được $10\text{ gam}$ kết tủa. Khối lượng sau phản ứng giảm $3{,}4\text{ gam}$ so với khối lượng dung dịch nước vôi trong ban đầu. Giá trị của $m$ là bao nhiêu?
\loigiai{}
\end{ex}

\begin{ex}
Có bao nhiêu amine bậc 1 trong các amine sau?\\
(1). $\ce{CH3CH2CH2NH2}$. \quad (2). $\ce{CH3CH(NH2)CH3}$. \quad (3). $\ce{CH3NHCH2CH3}$. \quad (4). $\ce{(CH3)3N}$.
\loigiai{}
\end{ex}

\begin{ex}
Một ester X mạch hở, có công thức phân tử $\ce{C4H8O2}$. Có bao nhiêu đồng phân cấu tạo ester của X?
\loigiai{}
\end{ex}

\begin{ex}
Có bao nhiêu đồng phân amine bậc một ứng với công thức phân tử $\ce{C4H11N}$?
\loigiai{}
\end{ex}

\begin{ex}
Cơ thể người mã hoá được mấy loại amino acid có trong bảng sau để tổng hợp protein cho cơ thể?
\begin{center}
	\includegraphics[width=14.0cm]{San_Pham/images_crop_485/fig_cau5_amino_acids_clean.png}
\end{center}
\loigiai{}
\end{ex}

\begin{ex}
Thuỷ phân hết $m\text{ gam}$ tetrapeptide Ala-Ala-Ala-Ala (mạch hở) thu được hỗn hợp gồm $28{,}24\text{ gam}$ Ala, $32\text{ gam}$ Ala-Ala và $27{,}27\text{ gam}$ Ala-Ala-Ala. Giá trị của $m$ là bao nhiêu? \textit{(Làm tròn kết quả 1 lần cuối cùng đến hàng đơn vị)}.
\loigiai{}
\end{ex}

\begin{ex}
Saccharose tham gia được bao nhiêu phản ứng nào sau đây?
\begin{enumerate}[(a)]
	\item Phản ứng với thuốc thử Tollens.
	\item Phản ứng với nước bromine.
	\item Phản ứng với $\ce{Cu(OH)2}$ tạo dung dịch màu xanh lam.
	\item Phản ứng thuỷ phân trong môi trường kiềm.
\end{enumerate}
\loigiai{}
\end{ex}

\begin{ex}
Protein tham gia được loại phản ứng dưới đây? Viết theo số thứ tự tăng dần (1234...)
\begin{enumerate}[(1)]
	\item Phản ứng khử thành alcohol.
	\item Phản ứng màu với $\ce{Cu(OH)2}$.
	\item Phản ứng màu với $\ce{HNO3}$.
	\item Phản ứng thuỷ phân.
\end{enumerate}
\loigiai{}
\end{ex}

\vspace{0.4cm}
\noindent
\textbf{B. PHẦN IV - TỰ LUẬN (3 ĐIỂM)}

\setcounter{ex}{0}

\begin{ex}
\textbf{(1,5 điểm)}
\begin{enumerate}[a)]
	\item Cho hợp chất hữu cơ X có công thức phân tử $\ce{C3H6O2}$. Hãy viết các đồng phân cấu tạo ester và gọi tên chúng.
	\item Viết công thức khung phân tử của stearic acid và oleic acid, biết oleic acid là acid béo omega-9, có liên kết đôi $\ce{C=C}$ ở dạng cis. Dự đoán nhiệt độ nóng chảy của oleic acid và stearic acid. Giải thích.
\end{enumerate}
\loigiai{}
\end{ex}

\begin{ex}
\textbf{(0,75 điểm)}
\begin{enumerate}[a)]
	\item Cho các lọ đựng các dung dịch sau bị mất nhãn: glucose, fructose và tinh bột. Trình bày cách nhận biết chúng.
	\item X là hỗn hợp gồm glucose, fructose và tinh bột. Phân tích thành phần nguyên tố trong X được $\%C = 40{,}82\%$; $\%H = 6{,}57\%$ (theo khối lượng). Phần trăm khối lượng tinh bột trong X là bao nhiêu?
\end{enumerate}
\loigiai{}
\end{ex}

\begin{ex}
\textbf{(0,75 điểm)} Cho $0{,}02\text{ mol}$ $\alpha$-amino acid X tác dụng vừa đủ với dung dịch chứa $0{,}04\text{ mol}$ $\ce{NaOH}$. Mặt khác $0{,}02\text{ mol}$ X tác dụng vừa đủ với $0{,}02\text{ mol}$ $\ce{HCl}$, thu được $3{,}67\text{ gam}$ muối.
\begin{enumerate}[a)]
	\item Nhúng Quỳ tím vào dung dịch X thì quỳ tím chuyển sang màu gì? Giải thích.
	\item Xác định công thức cấu tạo và gọi tên X, Viết kí hiệu của X.
\end{enumerate}
\loigiai{}
\end{ex}

\begin{center}
\textbf{--------- HẾT ---------}
\end{center}

\end{document}
